\honest{The Parable of the Dagger}{Eliezer Yudkowsky}{2008}

This is a story. A court jester who dabbles in logic gives the king two boxes, one with gold and one with an angry frog. Each box bears an inscription about the boxes and about which inscription is false, and the jester adds that ``one, and only one, of the inscriptions is true.'' The king opens the wrong box and is savaged by the frog. \nb{The puzzle has one answer: whichever inscription is true, the gold is in the second box.} The jester starts to explain, and the king has him thrown in the dungeon.

The next day the king shows the jester two boxes, one with the key to his chains and one with a dagger for his heart. The first box reads ``Either both inscriptions are true, or both inscriptions are false.'' The second reads ``This box contains the key.'' The king says nothing about whether the inscriptions are true.

The jester reasons that the second inscription is true whether the first is true or false, so the key must be in the second box. He opens it and finds the dagger. ``It's logically impossible!'' he cries. \nb{His reasoning is valid, but it assumes that each inscription is either true or false. With the dagger in the second box, the first inscription can be neither. What is impossible is the assumption, not the dagger's place.}

``It is entirely possible,'' the king replies. ``I merely wrote those inscriptions on two boxes, and then I put the dagger in the second one.'' I do not explain further.

I note that the story is adapted from Raymond Smullyan. \nb{It is problem 70 of Smullyan's \textsc{What Is the Name of This Book?} (1978), with Portia's caskets. Smullyan's solution gives both explanations: the suitor had no information about the truth of the inscriptions, and he assumed each was either true or false. ``I can take any number of caskets that I please,'' Smullyan writes, ``and put an object in one of them and then write any inscriptions at all on the lids.'' The first half of the post, where a promise about the inscriptions makes the puzzle solvable, is Yudkowsky's addition, and it shows the contrast Smullyan draws.}
